\documentclass[10pt,landscape,a4paper]{article}
\usepackage[utf8]{inputenc}
\usepackage[T1]{fontenc}
\usepackage{geometry}
\usepackage{multicol}
\usepackage{amsmath,amssymb}
\usepackage{enumitem}
\usepackage{graphicx}
\usepackage{xcolor}
\usepackage{listings}
\usepackage{booktabs}
\usepackage{titlesec}
\usepackage{fancyhdr}
\usepackage{hyperref}
\usepackage{CJKutf8}

% Page layout
\geometry{top=0.8cm, bottom=0.8cm, left=0.8cm, right=0.8cm}
\pagestyle{empty}
\setlength{\parindent}{0pt}
\setlength{\parskip}{2pt}
\setlength{\columnsep}{8pt}

% Section formatting
\titleformat{\section}{\normalsize\bfseries\color{blue!70!black}}{}{0em}{}[\titlerule]
\titleformat{\subsection}{\small\bfseries\color{blue!50!black}}{}{0em}{}
\titlespacing{\section}{0pt}{4pt}{2pt}
\titlespacing{\subsection}{0pt}{3pt}{1pt}

% Code listing style
\lstset{
    language=Python,
    basicstyle=\ttfamily\scriptsize,
    keywordstyle=\color{blue!70!black}\bfseries,
    stringstyle=\color{red!60!black},
    commentstyle=\color{green!50!black}\itshape,
    numberstyle=\tiny\color{gray},
    backgroundcolor=\color{gray!8},
    frame=single,
    framerule=0.3pt,
    rulecolor=\color{gray!50},
    breaklines=true,
    breakatwhitespace=false,
    tabsize=2,
    showstringspaces=false,
    aboveskip=3pt,
    belowskip=3pt,
    xleftmargin=2pt,
    xrightmargin=2pt,
}

% Custom box for key concepts
\newcommand{\keybox}[1]{%
    \begin{center}
    \fcolorbox{blue!50!black}{blue!5}{%
        \parbox{0.95\linewidth}{\small #1}%
    }
    \end{center}
}

% Chinese hint style (gray, smaller)
\newcommand{\zhint}[1]{{\scriptsize\color{gray!70!black}#1}}

\begin{document}
\begin{CJK}{UTF8}{gbsn}

\begin{center}
    {\large\bfseries IERG4320/IEMS5726 --- Chapter 2: Basic Statistical Test --- Cheat Sheet}\\[2pt]
    {\small Prof.\ YAN Wenjing, Dr.\ Danny Yip \quad|\quad Department of Information Engineering, CUHK}
\end{center}

\vspace{2pt}
\begin{multicols}{3}

% ============================================================
\section{Key Terminology \zhint{(基础概念)}}
% ============================================================

\begin{itemize}[leftmargin=*, nosep, itemsep=1pt]
    \item \textbf{Population} \zhint{(总体)}: all data you want to study\\
    \zhint{比如"全校所有学生"}
    \item \textbf{Sample} \zhint{(样本)}: a portion you actually collect\\
    \zhint{比如"随机抽了100个学生"}
    \item \textbf{Variable} \zhint{(变量)}: anything you can measure\\
    \zhint{比如身高、成绩、年龄}
    \item \textbf{Descriptive stats} \zhint{(描述统计)}: describe what data looks like\\
    \zhint{比如"平均分是75"}
    \item \textbf{Inferential stats} \zhint{(推断统计)}: use sample to predict population\\
    \zhint{用100人的数据推测全校情况}
\end{itemize}

\subsection{Central Tendency \zhint{(集中趋势)}}
\begin{itemize}[leftmargin=*, nosep, itemsep=1pt]
    \item \textbf{Mean} \zhint{(均值/平均数)}: $\bar{x} = \frac{1}{n}\sum_{i=1}^{n} x_i$\\
    \zhint{所有数加起来除以个数}
    \item \textbf{Median} \zhint{(中位数)}: middle value when sorted\\
    \zhint{数据排好序后最中间那个}
    \item \textbf{Mode} \zhint{(众数)}: most frequent value\\
    \zhint{出现次数最多的那个值}
\end{itemize}

\subsection{Spread \zhint{(离散程度)}}
\begin{itemize}[leftmargin=*, nosep, itemsep=1pt]
    \item \textbf{Range} \zhint{(范围)}: $\max - \min$
    \item \textbf{Variance} \zhint{(方差)}: $\sigma^2 = \frac{1}{n}\sum(x_i - \bar{x})^2$\\
    \zhint{每个值离均值有多远，取平方再平均}
    \item \textbf{Std Dev} \zhint{(标准差)}: $\sigma = \sqrt{\sigma^2}$\\
    \zhint{方差开根号，更直观看数据有多分散}
\end{itemize}

% ============================================================
\section{Hypothesis Testing \zhint{(假设检验)}}
% ============================================================

\zhint{核心问题：观察到的结果到底是真的有意义，还是碰巧？}

\keybox{
\textbf{Goal} \zhint{(目的)}: Rule out ``luck'' as the explanation.\\
\zhint{排除"纯属巧合"这种可能性。}
}

\subsection{6 Steps \zhint{(六步走)}}
\begin{enumerate}[leftmargin=*, nosep, itemsep=1pt]
    \item \textbf{State Hypotheses} \zhint{(提出假设)}
    \begin{itemize}[nosep, itemsep=0pt]
        \item $H_0$ \zhint{(零假设)}: no effect \zhint{(没效果)}
        \item $H_a$ \zhint{(备择假设)}: effect exists \zhint{(有效果)}
    \end{itemize}
    \item \textbf{Choose $\alpha$} \zhint{(选显著性水平)}: 0.05, 0.01, or 0.1\\
    \zhint{$\alpha=0.05$意思是：巧合概率$<$5\%时，我就认为不是巧合}
    \item \textbf{Pick a test} \zhint{(选检验方法)}: Z, t, $\chi^2$, etc.
    \item \textbf{Calculate test statistic} \zhint{(算统计量)}
    \item \textbf{Find critical region} \zhint{(确定拒绝域)}
    \item \textbf{Decision} \zhint{(做决定)}: reject $H_0$ or not
\end{enumerate}

\subsection{Critical Regions \zhint{(拒绝域)}}
\textbf{One-tailed} \zhint{(单尾，只看一个方向)}:
\begin{itemize}[leftmargin=*, nosep, itemsep=0pt]
    \item Left: reject $H_0$ if $z < -z_\alpha$
    \item Right: reject $H_0$ if $z > z_\alpha$
\end{itemize}
\textbf{Two-tailed} \zhint{(双尾，两边都看)}: reject if $|z| > z_{\alpha/2}$\\
\zhint{比如检验硬币公不公平，正面太多或太少都算有问题}

\subsection{p-value \zhint{(p值)}}
\keybox{
$p$-value = probability of seeing this result (or more extreme) when $H_0$ is true.\\
\zhint{p值 = 假如真的没效果，出现当前结果的概率有多大}\\[2pt]
\textbf{Rule}: $p < \alpha \Rightarrow$ reject $H_0$ \zhint{(拒绝零假设 = 有显著效果)}\\
\zhint{p值越小，越能说明"不是巧合"}
}

\textbf{Example} \zhint{(例子)}: Test if coin is fair ($H_0: p=0.5$), flip 100 times, get 60 heads:\\
$p = 2\,P(X \geq 60) = 0.0574 > 0.05$\\
$\Rightarrow$ Cannot reject $H_0$ \zhint{(不能说硬币不公平)}

% ============================================================
\section{Fisher's Exact Test \zhint{(Fisher精确检验)}}
% ============================================================

\zhint{适用场景：2$\times$2表格、小样本、变量只有两类（是/否）}

\begin{itemize}[leftmargin=*, nosep, itemsep=1pt]
    \item For \textbf{2$\times$2 contingency tables} \zhint{(列联表)}
    \item Tests if two binary variables are independent \zhint{(两个变量是否有关)}
    \item Computes \textbf{exact} probability \zhint{(精确计算，不是近似)}
\end{itemize}

\subsection{Contingency Table \zhint{(列联表)}}
\begin{center}
\begin{tabular}{l|cc|c}
\toprule
 & Group I & Group II & Total \\
\midrule
Category A & $a$ & $b$ & $a+b$ \\
Category B & $c$ & $d$ & $c+d$ \\
\midrule
Total & $a+c$ & $b+d$ & $n$ \\
\bottomrule
\end{tabular}
\end{center}

\subsection{Example \zhint{(例子：骰子公不公平？)}}
\begin{center}
\begin{tabular}{l|cc}
\toprule
Rolling 2 dice & Exper. & Theory \\
\midrule
Sum $\geq$ 10 & 15 & 6 \\
Sum $<$ 10 & 57 & 30 \\
\bottomrule
\end{tabular}
\end{center}

\begin{lstlisting}
import scipy.stats as stats
oddsratio, pvalue = stats.fisher_exact(
    [[15, 57], [6, 30]], "greater")
print(oddsratio, " ", pvalue)
# 1.3158  0.4051
\end{lstlisting}
$p = 0.405 \gg 0.05 \Rightarrow$ Cannot reject $H_0$\\
\zhint{p值远大于0.05，没有证据说骰子不公平}

\subsection{Pros \& Cons \zhint{(优缺点)}}
\zhint{优点：小样本精确；缺点：只能处理2$\times$2表格，大数据慢}

% ============================================================
\section{Chi-squared Test \zhint{(卡方检验, $\chi^2$)}}
% ============================================================

\zhint{适用场景：多类别频率比较，看实际频率和预期频率是否有显著差异}

\subsection{Test Statistic \zhint{(公式)}}
\[
\chi^2 = \sum_{i=1}^{k} \frac{(O_i - E_i)^2}{E_i}
\]
\begin{itemize}[leftmargin=*, nosep, itemsep=0pt]
    \item $O_i$: observed count \zhint{(实际观察到的次数)}
    \item $E_i = N p_i$: expected count \zhint{(理论上期望的次数)}
    \item $df = k - 1$ \zhint{(自由度 = 类别数 $-$ 1)}
\end{itemize}
\zhint{核心思想：$\chi^2$越大，说明实际和预期差得越多}

\subsection{Chi-square Distribution \zhint{(卡方分布)}}
If $X_1, \ldots, X_k \sim N(0,1)$ independent, then:
\[
Y = X_1^2 + X_2^2 + \cdots + X_k^2 \sim \chi^2(k)
\]
$E[Y]=k$, \; $\mathrm{Var}(Y)=2k$

\subsection{Example \zhint{(例子：骰子有没有偏？)}}
\zhint{骰子掷60次，各面出现：5, 8, 9, 8, 10, 20次。公平骰子每面应该10次。}

\begin{center}
\small
\begin{tabular}{c|cccccc|c}
\toprule
$i$ & 1 & 2 & 3 & 4 & 5 & 6 & $\Sigma$ \\
\midrule
$O_i$ & 5 & 8 & 9 & 8 & 10 & 20 & 60 \\
$E_i$ & 10 & 10 & 10 & 10 & 10 & 10 & 60 \\
$\frac{(O-E)^2}{E}$ & 2.5 & 0.4 & 0.1 & 0.4 & 0 & 10 & \textbf{13.4} \\
\bottomrule
\end{tabular}
\end{center}
\zhint{注意：6出现了20次，远超预期，是主要偏差来源}

\begin{lstlisting}
from scipy.stats import chisquare
chi2, pvalue = chisquare([5,8,9,8,10,20])
print(chi2, " ", pvalue)
# 13.4  0.01991
\end{lstlisting}

\textbf{Decision} ($df=5$, $\chi^2 = 13.4$):
\begin{itemize}[leftmargin=*, nosep, itemsep=0pt]
    \item 95\%: $13.4 > 11.07 \Rightarrow$ \textbf{Reject $H_0$} \zhint{(骰子不公平!)}
    \item 99\%: $13.4 < 15.09 \Rightarrow$ Cannot reject \zhint{(没法下结论)}
\end{itemize}
\zhint{同样的数据，标准不同结论就不同，这就是$\alpha$的意义}

% ============================================================
\section{Wilcoxon Rank-Sum Test \zhint{(秩和检验)}}
\small{(Mann--Whitney U Test)}
% ============================================================

\zhint{适用场景：两组独立数据的比较，数据不需要正态分布（t检验的替代品）}

\subsection{Key Idea \zhint{(核心思想)}}
\begin{itemize}[leftmargin=*, nosep, itemsep=0pt]
    \item $H_0$: $P(X > Y) = P(Y > X)$ \zhint{(两组没差别)}
    \item Both groups must be independent \zhint{(两组互相独立)}
\end{itemize}
\zhint{如果两组没差别，混合排名后两组的排名应该差不多}

\subsection{Steps \zhint{(步骤)}}
\begin{enumerate}[leftmargin=*, nosep, itemsep=1pt]
    \item Mix both groups, rank ascending \zhint{(混一起从小到大排名)}\\
    \zhint{相同值取排名的平均}
    \item $T$ = sum of ranks of one group \zhint{(一组的排名加起来)}
    \item $U = T - \frac{n(n+1)}{2}$ \zhint{(算U统计量)}
    \item $\mu_U = \frac{n_1 n_2}{2}$, \; $\sigma_U = \sqrt{\frac{n_1 n_2 (n_1+n_2+1)}{12}}$
    \item $z = \frac{U - \mu_U}{\sigma_U}$ $\to$ get $p$-value
\end{enumerate}

\subsection{Example}
$X$ (15 values), $Y$ (14 values):\\
$U_X = 152$, $\mu_U = 105$, $\sigma_U = 56.13$\\
$z = 0.837$, $p = 0.040$\\
\zhint{在99\%水平下不能拒绝$H_0$，没法说两组有差别}

\begin{lstlisting}
from scipy.stats import ranksums, mannwhitneyu
import numpy as np
X = np.array([3.99, 3.79, 3.60, 3.73, 3.21,
    3.60, 4.08, 3.61, 3.83, 3.31,
    4.13, 3.26, 3.54, 3.51, 2.71])
Y = np.array([3.18, 2.84, 2.90, 3.27, 3.85,
    3.52, 3.23, 2.76, 3.60, 3.75,
    3.59, 3.63, 2.38, 2.34])

# Method 1: ranksums
U, pvalue = ranksums(X, Y)
print(U, pvalue)  # 2.05  0.0402

# Method 2: mannwhitneyu
stat, pvalue = mannwhitneyu(X, Y)
print(stat, pvalue)  # 152.0  0.0423
\end{lstlisting}
\zhint{两个函数结果略不同但结论一致}

% ============================================================
\section{Wilcoxon Signed-Rank Test \zhint{(符号秩检验)}}
% ============================================================

\zhint{适用场景：配对数据（同一组人的前后对比），数据不需要正态分布（配对t检验的替代品）}

\subsection{Key Idea \zhint{(核心思想)}}
\begin{itemize}[leftmargin=*, nosep, itemsep=0pt]
    \item $H_0$: paired samples come from same distribution\\
    \zhint{两组配对数据来自相同分布（处理没效果）}
    \item Data must be paired \zhint{(数据必须配对，比如同一个人的前/后)}
\end{itemize}
\zhint{如果处理没效果，正的差值和负的差值应该差不多抵消}

\subsection{Steps \zhint{(步骤)}}
\begin{enumerate}[leftmargin=*, nosep, itemsep=1pt]
    \item $d_i = x_i - y_i$ \zhint{(算每对的差值)}
    \item Drop pairs where $d_i = 0$ \zhint{(差为0的扔掉)}
    \item Rank by $|d_i|$ ascending \zhint{(按差值绝对值排名)}
    \item $W = \sum$ signed ranks \zhint{(加上正负号的排名求和)}
    \item $\sigma_W = \sqrt{\frac{n(n+1)(2n+1)}{6}}$
    \item $t = W / \sigma_W$, use $t$-dist ($df = n-1$)
\end{enumerate}

\subsection{Example}
\begin{center}
\scriptsize
\begin{tabular}{cccc}
\toprule
$X$ & $Y$ & $X-Y$ & Rank \\
\midrule
6.1 & 5.2 & 0.9 & 3.5 \\
7.0 & 7.9 & $-0.9$ & 3.5 \\
8.2 & 3.9 & 4.3 & 10 \\
7.6 & 4.7 & 2.9 & 7 \\
6.5 & 5.3 & 1.2 & 5 \\
8.4 & 5.4 & 3.0 & 8 \\
6.9 & 4.2 & 2.7 & 6 \\
6.7 & 6.1 & 0.6 & 2 \\
7.4 & 3.8 & 3.6 & 9 \\
5.8 & 6.3 & $-0.5$ & 1 \\
\bottomrule
\end{tabular}
\end{center}

$W = 46$, $\sigma_W = 19.62$, $df = 9$, $p = 0.0438$\\
\zhint{在99\%水平下不能拒绝$H_0$}

\begin{lstlisting}
from scipy.stats import wilcoxon
import numpy as np
X = np.array([6.1, 7.0, 8.2, 7.6, 6.5,
              8.4, 6.9, 6.7, 7.4, 5.8])
Y = np.array([5.2, 7.9, 3.9, 4.7, 5.3,
              5.4, 4.2, 6.1, 3.8, 6.3])
T, pvalue = wilcoxon(X, Y)
print(T, pvalue)  # 5.0  0.01953
\end{lstlisting}
\zhint{注意：Python用T统计量（正/负秩和的较小值），和手算W不同，但结论一致}

% ============================================================
\section{Quick Reference \zhint{(速查表：用哪个检验？)}}
% ============================================================

\begin{center}
\small
\begin{tabular}{lp{4.2cm}}
\toprule
\textbf{Test} & \textbf{When to Use} \zhint{(什么时候用)} \\
\midrule
Fisher's Exact & 2$\times$2 table, small $n$, binary \zhint{(小样本，是/否变量)} \\
$\chi^2$ Test & Multi-category frequency \zhint{(多类别频率比较)} \\
Rank-Sum (U) & 2 independent groups, non-normal \zhint{(两组独立数据)} \\
Signed-Rank & Paired data, non-normal \zhint{(配对数据)} \\
\bottomrule
\end{tabular}
\end{center}

\keybox{
\textbf{Core logic} \zhint{(核心套路)}:\\
Compute statistic $\to$ get $p$-value $\to$ if $p < \alpha$, reject $H_0$\\
\zhint{算统计量 $\to$ 得到p值 $\to$ p值$<\alpha$就拒绝零假设 $\to$ 结果不是巧合!}
}

\subsection{Other Tests \zhint{(其他检验，了解即可)}}
\begin{itemize}[leftmargin=*, nosep, itemsep=0pt]
    \item \textbf{t-test} \zhint{(t检验)}: compare two groups (need normal data)
    \item \textbf{ANOVA} \zhint{(方差分析)}: compare $>$2 groups
    \item \textbf{Correlation} \zhint{(相关分析)}: are two variables linearly related?
    \item \textbf{Regression} \zhint{(回归)}: predict Y from X
\end{itemize}

\end{multicols}
\end{CJK}
\end{document}
